%Structure d'une entrée de glossaire
\newglossaryentry{bibliothèque numérique}{
		name=bibliothèque numérique,
		sort=bibliothèque numérique,
		description={Les bibliothèques numériques sont des organisations qui offrent des ressources, y compris en personnel, pour sélectionner, structurer, offrir un accès intellectuel, interpréter, distribuer et conserver l'intégrité de documents sous forme numérique. Une bibliothèque numérique garantit également un accès sur la durée aux oeuvres électroniques dans le but d'être aisément et à moindre coût disponibles à un ou plusieurs publics spécifiques\footcite{claerrManuelConstitutionBibliotheques2013a}.}
}

\newglossaryentry{Bibliostratus}{
		name=Bibliostratus,
		sort=Bibliostratus,
		description={Outils logiciels d'alignement et de nettoyer de données massives, utilisés pour transformer les formats locaux et lier les autorités de la bibliothèque aux référentiels nationaux.}
}

\newglossaryentry{OpenRefine}{
		name=OpenRefine,
		sort=OpenRefine,
		description={Outils logiciels d'alignement et de nettoyage de données massives, utilisés pour transformer les formats locaux et lier les autorités de la bibliothèque aux référentiels nationaux.}
}

\newglossaryentry{Dublin Core}{
		name=Dublin Core,
		sort=Dublin Core,
		description={Schéma de métadonnées génériques comprenant quinze éléments de base.C'est un schéma de métadonnées générique et standardisé conçu pour décrire des ressources numériques ou physiques.}
}

\newglossaryentry{IdRef}{
		name=IdRef,
		sort=IdRef,
		description={Une application Web développée et maintenue par l'ABES. IdRef permet à des utilisateurs et à des applications tierces d'interroger, de consulter, de créer et d'enrichir des autorités.}
}

\newglossaryentry{IFLA-LRM}{
		name=IFLA-LRM,
		sort=IFLA-LRM,
		description={Un modèle conceptuel de haut niveau publié en 2017 par la Fédération internationale des associations et institutions de bibliothèques (IFLA). Il décrit la structure logique de l'univers bibliographique en définissant des entités, des attributs et des relations.}
}

\newglossaryentry{IIIF}{
		name=IIIF,
		sort=IIIF,
		description={Un ensemble de règles techniques et une communauté internationale qui permettent de partager et de manipuler facilement des images en haute résolution sur le Web.}
}

\newglossaryentry{OAI-PMH}{
		name=OAI-PMH,
		sort=OAI-PMH,
		description={Protocole informatique standardisé qui permet d'échanger des métadonnées (des données décrivant des documents) sur Internet.}
}

\newglossaryentry{RAMEAU}{
		name=RAMEAU,
		sort=RAMEAU,
		description={Langage d'indexation matière utilisé, en France, par la Bibliothèque nationale de France, les bibliothèques universitaires, de nombreuses bibliothèques de lecture publique ou de recherche ainsi que plusieurs organismes privés.}
}

\newglossaryentry{RDF}{
		name=RDF,
		sort=RDF,
		description={Modèle de graphe destiné à décrire formellement les ressources Web et leurs métadonnées, afin de permettre le traitement automatique de ces descriptions.}
}

\newglossaryentry{SKOS}{
		name=SKOS,
		sort=SKOS,
		description={Recommandation du W3C publiée le 18 août 2009 pour représenter des thésaurus documentaires, classifications ou d'autres types de vocabulaires contrôlés ou de langages documentaires. S'appuyant sur le modèle de données RDF, son principal objectif est de permettre la publication facile de vocabulaires structurés pour leur utilisation dans le cadre du Web sémantique.}
}

\newglossaryentry{Web sémantique}{
		name=Web sémantique,
		sort=Web sémantique,
		description={Un modèle qui permet aux données d'être partagées et réutilisées entre plusieurs applications, entreprises et groupes d'utilisateurs.}
}

\newglossaryentry{Web des données}{
		name=Web des données,
		sort=Web des données,
		description={Le Web des données \textit{linked data} en anglais est une initiative du World Wide Web Consortium (W3C) visant à favoriser la publication de données structurées sur le Web, non pas sous la forme de silos de données isolés les uns des autres, mais en les reliant entre elles pour constituer un réseau global d'informations.}
}

\newglossaryentry{XTF}{
		name=XTF,
		sort=XTF,
		description={Le terme XTF désigne principalement une plateforme open source de gestion, d'indexation et de recherche pour les bibliothèques numériques.}
}